% Top-level content for the online blueprint.

\blueprintfrontmatter

\begin{abstract}
This document gives a proof of \textbf{Kummer's criterion}:
\[
  p \text{ is regular}
  \quad\Longleftrightarrow\quad
  \forall k,\ 1 \le k,\ 2k \le p - 3,\quad
  p \nmid \operatorname{num}(B_{2k})
\]
for every odd prime $p$.

The proof has four steps.  First, for a prime cyclotomic CM field $K$ and its
maximal real subfield $K^+$, extension of ideals gives an injective class-group
map
\[
  \Cl(\mathcal O_{K^+}) \longrightarrow \Cl(\mathcal O_K).
\]
This proves $h^+ \mid h$ and allows the relative class number
\[
  h^- := h/h^+
\]
to be used as an actual natural number.

Second, the relative class-number formula is reduced to a $p$-adic congruence
and then to the Bernoulli numerator criterion:
\[
  p \mid h^- \quad\Longleftrightarrow\quad
  \exists k,\ 1 \le k,\ 2k \le p-3,\quad
  p \mid \operatorname{num}(B_{2k}).
\]

Third, the real cyclotomic units form a subgroup $C^+$ of the full unit group
of $K^+$.  The Kummer logarithm matrix proves that Bernoulli numerator
non-divisibility forces $C^+$ to be $p$-saturated, hence
\[
  p \nmid [\mathcal O_{K^+}^{\times}:C^+].
\]
The prime-conductor cyclotomic-unit index theorem identifies the odd-primary
part of this index with $h^+$.

Finally, these ingredients show that $p \mid h$ is equivalent to divisibility
of one Bernoulli numerator in Kummer's range.  The definition of regular prime
then gives the stated criterion.

The Lean declaration corresponding to the final theorem is
\texttt{KummerCriterion}.

As an application of the criterion, a final chapter proves Carlitz's theorem
that there are infinitely many irregular primes, via the unrestricted Kummer
congruence for divided Bernoulli numbers.
\end{abstract}

\clearpage
\setcounter{page}{2}
\tableofcontents

\blueprintmainmatter
\chapter{Introduction}\label{ch:intro}
% Introduction: history, main results, notation, and outline.

\section{Overview}

Let $p$ be an odd prime, let $\zeta_p$ be a primitive $p$th root of unity,
and let $K=\Qzetap$ be the $p$th cyclotomic field.  Kummer
\cite{kummer1850} proved that Fermat's Last Theorem holds for the exponent
$p$ whenever $p$ is \emph{regular}, that is, whenever $p$ does not divide
the class number $h(K)$.  Regularity is a condition on an invariant that is
notoriously difficult to compute; Kummer's second fundamental insight is
that it can be detected by a finite, purely arithmetic test on Bernoulli
numbers.

This document is the blueprint of a formalisation, in Lean~4 and building
on mathlib, of \emph{Kummer's criterion}: an odd prime $p$ is regular if
and only if
\[
  p \nmid \operatorname{num}(B_{2k})
  \qquad\text{for all } k \text{ with } 1 \le k,\ 2k \le p-3,
\]
where $B_n$ denotes the $n$th Bernoulli number and $\operatorname{num}$ the
numerator of a rational number in lowest terms.  The final statement is
\cref{thm:kummer-criterion}; the formal counterpart is the Lean declaration
\texttt{KummerCriterion}.

The criterion makes regularity decidable in practice.  Running the
Bernoulli test on the primes below $100$ shows that the irregular primes in
that range are exactly $37$, $59$ and $67$; combined with the
\texttt{flt-regular} formalisation of Kummer's theorem on Fermat's Last
Theorem for regular prime exponents, this yields Fermat's Last Theorem for
every exponent $2 < n \le 100$ outside $\{37, 59, 67, 74\}$.

Irregular primes are not a finite accident.  Jensen \cite{jensen1915}
proved that there are infinitely many irregular primes (indeed, infinitely
many congruent to $3$ modulo $4$), and Carlitz \cite{carlitz1954} gave the
short argument formalised here, which combines the criterion with the
unrestricted Kummer congruence for divided Bernoulli numbers and the
archimedean growth of $|B_{2k}/2k|$.  The final statement is
\cref{thm:infinitely-many-irregular}; the formal counterpart is
\texttt{KummerCriterion.infinite\_not\_isRegularPrime}.

Standard references for all of this material are Washington
\cite{washington} and Ireland--Rosen \cite{ireland-rosen}.

\begin{definition}\label{def:regular-prime}
  \lean{IsRegularPrime}\leanok
An odd prime $p$ is \emph{regular} when $p$ does not divide the class
number of the $p$th cyclotomic field:
\[
  p \nmid h\bigl(\Qzetap\bigr).
\]
Otherwise $p$ is \emph{irregular}.  (In the formalisation the condition is
phrased as coprimality of $p$ with the cardinality of the class group of
$\Qzetap$, which is equivalent because $p$ is prime.)
\end{definition}

\section{Notation and conventions}

Throughout the blueprint:

\begin{itemize}
\item $p$ is an odd prime, $\zeta_p$ a primitive $p$th root of unity,
  $K = \Qzetap$, and $K^+$ the maximal real subfield of $K$
  (\cref{def:cm-real-subfield}).  We write $\mathcal O_K$, $\mathcal
  O_{K^+}$ for the rings of integers and $\Cl{\mathcal O_K}$,
  $\Cl{\mathcal O_{K^+}}$ for the ideal class groups.
\item $h = h(K)$, $h^+ = h(K^+)$ and $h^- = h/h^+$ are the class number,
  the plus class number, and the relative class number
  (\cref{def:class-number-h}, \cref{def:class-number-hplus},
  \cref{def:class-number-hminus}).
\item $B_n$ is the $n$th Bernoulli number, with the convention
  $B_1 = -\tfrac12$ (the convention of mathlib's \texttt{bernoulli}).
  Thus $B_n = 0$ for odd $n \ge 3$, and $B_n(X)$ denotes the $n$th
  Bernoulli polynomial.  For a rational number $q$ written in lowest terms
  we write $\operatorname{num}(q) \in \Z$ and $\operatorname{den}(q) \in
  \N$ for its numerator and denominator.
\item Dirichlet characters modulo $p$ (\cref{def:dirichlet-character}) are
  written $\chi$; the trivial character is $1$, and $\chi$ is \emph{even}
  or \emph{odd} according to $\chi(-1) = 1$ or $\chi(-1) = -1$.  The
  generalised Bernoulli numbers $B_{n,\chi}$ are defined in
  \cref{def:generalized-bernoulli}, the Teichm\"uller character $\omega$ in
  \cref{def:teichmuller}.  Gauss sums are taken with respect to the
  standard additive character $a \mapsto e^{2\pi i a/p}$ and written
  $\tau(\chi)$.
\item $\Q_p$ is the field of $p$-adic numbers and $\Z_p$ its ring of
  integers.  For rational numbers $x, y$ (or, more generally, elements of
  $\Q_p$) we write
  \[
    x \equiv y \pmod{p}
  \]
  to mean $x - y \in p\Z_p$.  In the formalisation this is rendered as the
  existence of $z \in \Z_p$ with $x - y = pz$ in $\Q_p$.
\item For a finite abelian group quotient we write $[E : H]$ for the index
  of a subgroup $H \le E$.
\end{itemize}

\section{Outline}

The proof of the criterion has four parts, one per chapter.

Chapter~\ref{ch:cm} proves that extension of ideals induces an
\emph{injective} map $\Cl{\mathcal O_{K^+}} \to \Cl{\mathcal O_K}$
(\cref{thm:class-group-map-injective}).  Consequently $h^+ \mid h$, and the
relative class number $h^- = h/h^+$ is a genuine natural number with
$h = h^+ h^-$.

Chapter~\ref{ch:hminus-formula} establishes the analytic formula
\[
  h^- = 2p \prod_{\chi \text{ odd}} \Bigl(-\tfrac12 B_{1,\chi^{-1}}\Bigr)
\]
(\cref{thm:hminus-formula}), then reduces it modulo $p$: indexing the odd
characters by powers of the Teichm\"uller character and applying the
congruence $B_{1,\omega^j} \equiv B_{j+1}/(j+1) \pmod p$ converts the
product into a product of classical divided Bernoulli numbers, giving
\[
  p \mid h^- \iff
  \exists k,\ 1 \le k,\ 2k \le p-3,\ p \mid \operatorname{num}(B_{2k})
\]
(\cref{thm:hminus-bernoulli-final}).

Chapter~\ref{ch:cyclotomic-units} treats the plus side.  The real
cyclotomic units form a subgroup $C^+$ of $(\mathcal O_{K^+})^\times$, and
the determinant of the \emph{Kummer logarithm matrix} --- a matrix of
$p$-adic logarithms of the cyclotomic-unit generators --- factors as a
product of Bernoulli numbers times a nonzero Vandermonde determinant
(\cref{thm:kummer-log-det-bernoulli}).  If no Bernoulli numerator in
Kummer's range is divisible by $p$, the determinant is nonzero and $C^+$ is
$p$-saturated in the full unit group, whence
$p \nmid [(\mathcal O_{K^+})^\times : C^+]$
(\cref{thm:kummer-log-det-saturation} and
\cref{thm:not-dvd-index-of-pSaturated}).
Sinnott's prime-conductor index theorem \cite{sinnott1978} identifies the
$p$-divisibility of this index with that of $h^+$
(\cref{thm:sinnott-index-pprimary}).

Chapter~\ref{ch:kummer-final} assembles the criterion.  The plus-side
results yield the implication $p \mid h^+ \Rightarrow p \mid h^-$
(\cref{thm:hplus-to-hminus-units}); together with $h = h^+h^-$ this reduces
$p \mid h$ to $p \mid h^-$, and the minus criterion finishes the proof
(\cref{thm:kummer-criterion}).

Chapter~\ref{ch:irregular-primes} contains the application to the
infinitude of irregular primes, following Carlitz \cite{carlitz1954}: a
finite set $S$ of primes is escaped by manufacturing an even index $M$,
divisible by $q - 1$ for every prime $q \in S$, with $|B_M/M| > 1$; any
prime in the numerator of $B_M/M$ is then irregular and lies outside $S$.
The key arithmetic input is the unrestricted Kummer congruence
$B_m/m \equiv B_n/n \pmod p$ for even $m \equiv n \not\equiv 0
\pmod{p-1}$, proved by the elementary Voronoi route
(\cref{thm:kummer-congruence-full}).


\pagebreak
\chapter{Class numbers of the cyclotomic CM field}\label{ch:cm}
% CM splitting: the injection Cl(O_{K^+}) -> Cl(O_K) and h = h^+ h^-.

In this chapter $K$ denotes the $p$th cyclotomic field for an odd prime
$p$.  The goal is the factorisation $h = h^+ h^-$ of the class number: we
prove that the natural map on class groups induced by the inclusion
$\mathcal O_{K^+} \subseteq \mathcal O_K$ is injective, so that $h^+ \mid
h$ and the relative class number $h^- = h/h^+$ makes sense as a natural
number.  The key step is a descent statement for principal ideals
(\cref{thm:is-principal-after-extension}): an ideal of $\mathcal O_{K^+}$
that becomes principal in $\mathcal O_K$ was principal to begin with.

\section{The real subfield and the class numbers}

\begin{definition}\label{def:cm-real-subfield}
  \lean{NumberField.maximalRealSubfield}\leanok
For a number field $K$, let $K^+$ denote its maximal real subfield, the
largest subfield of $K$ all of whose complex embeddings are real.  When $K$
is a CM field, $K^+$ is the fixed field of complex conjugation, and
$[K : K^+] = 2$.  For $K = \Qzetap$ this is the usual maximal real
subfield $\Q(\zeta_p + \zeta_p^{-1})$.
\end{definition}

\begin{definition}\label{def:class-number-h}
  \lean{KummerCriterion.h}\leanok
For a number field $K$, the class number $h(K)$ is the cardinality of the
ideal class group of its ring of integers:
\[
  h(K)=\#\Cl{\mathcal O_K}.
\]
\end{definition}

\begin{definition}\label{def:class-number-hplus}
  \lean{KummerCriterion.hPlus}\leanok
  \uses{def:cm-real-subfield, def:class-number-h}
For a CM field $K$, the plus class number is the class number of the
maximal real subfield:
\[
  h^+(K)=h(K^+).
\]
\end{definition}

\section{Descent of principal ideals}

\begin{theorem}\label{thm:ring-of-integers-faithfully-flat}
  \lean{KummerCriterion.ringOfIntegers_faithfullyFlat_maximalRealSubfield}\leanok
  \uses{def:cm-real-subfield}
For every number field $L$, the ring $\mathcal O_L$ is faithfully flat over
$\mathcal O_{L^+}$.
\end{theorem}
\begin{proof}\leanok
$\mathcal O_L$ is a finitely generated torsion-free module over the
Dedekind domain $\mathcal O_{L^+}$, hence projective, hence flat.  For
faithfulness we use the criterion in terms of prime spectra: it suffices
that $\operatorname{Spec}\mathcal O_L \to \operatorname{Spec}\mathcal
O_{L^+}$ be surjective.  Given a prime $\mathfrak p$ of $\mathcal O_{L^+}$,
the extension $\mathcal O_{L^+} \subseteq \mathcal O_L$ is integral, so
lying-over produces a prime of $\mathcal O_L$ contracting to $\mathfrak p$.
\end{proof}

\begin{theorem}\label{thm:map-comap-ring-of-integers}
  \lean{KummerCriterion.map_comap_eq_ringOfIntegers}\leanok
  \uses{thm:ring-of-integers-faithfully-flat}
If $J$ is an ideal of $\mathcal O_{L^+}$, then extension to $\mathcal O_L$
followed by contraction is the identity:
\[
  (J\mathcal O_L)\cap\mathcal O_{L^+}=J.
\]
\end{theorem}
\begin{proof}\leanok
For a faithfully flat ring map $A \to B$, every ideal $J \subseteq A$
satisfies $JB \cap A = J$: the inclusion $J \subseteq JB \cap A$ is clear,
and the quotient map $A/J \to B/JB$ is injective by faithful flatness of
$B/JB$ over $A/J$.  Apply this with
\cref{thm:ring-of-integers-faithfully-flat}.
\end{proof}

\begin{theorem}\label{thm:is-principal-map-span-singleton}
  \lean{KummerCriterion.isPrincipal_of_map_eq_span_singleton_of_mem}\leanok
  \uses{thm:map-comap-ring-of-integers}
Let $I$ be an ideal of $\mathcal O_{K^+}$.  If $I\mathcal O_K$ is generated
by the image of an element $b_0\in\mathcal O_{K^+}$, then $I$ is principal,
generated by $b_0$.
\end{theorem}
\begin{proof}\leanok
Contract the equality $I\mathcal O_K=(b_0)$ back to $\mathcal O_{K^+}$.
The left-hand side contracts to $I$ by
\cref{thm:map-comap-ring-of-integers}.  For the right-hand side, the
contraction of $b_0\mathcal O_K$ contains $b_0\mathcal O_{K^+}$, and
conversely any element of $b_0 \mathcal O_K \cap \mathcal O_{K^+}$ is of
the form $b_0 c$ with $c \in K \cap \mathcal O_K = \mathcal O_{K^+}$
(using $b_0 \ne 0$; if $b_0 = 0$ then $I = 0$ is principal anyway).  Hence
$I=(b_0)$.
\end{proof}

\begin{theorem}\label{thm:ideal-map-conj}
  \lean{KummerCriterion.ideal_map_conj_eq}\leanok
  \uses{def:cm-real-subfield}
If $I$ is an ideal of $\mathcal O_{K^+}$, then the extended ideal
$I\mathcal O_K$ is fixed by complex conjugation.
\end{theorem}
\begin{proof}\leanok
The extension $I\mathcal O_K$ is generated by elements coming from $K^+$,
and those elements are fixed by conjugation.  Conjugation is a ring
automorphism of $\mathcal O_K$, so it maps $I\mathcal O_K$ onto the ideal
generated by the conjugates of the generators, which is $I\mathcal O_K$
itself.
\end{proof}

\begin{theorem}\label{thm:conj-generator-associated}
  \lean{KummerCriterion.conj_generator_associated}\leanok
  \uses{thm:ideal-map-conj}
Suppose $I \ne 0$ and $I\mathcal O_K=(a)$.  Then, by
\cref{thm:ideal-map-conj}, also $I\mathcal O_K=(\overline a)$, and the two
generators are associated: there is a unit $u\in\mathcal O_K^\times$ with
\[
  \overline a = u a.
\]
\end{theorem}
\begin{proof}\leanok
Two generators of the same nonzero principal ideal in a domain differ by a
unit: from $(a) = (\overline a)$ we get $\overline a = ua$ and $a = v
\overline a$ for some $u, v \in \mathcal O_K$, whence $a = vua$ and, since
$a \ne 0$, $vu = 1$.
\end{proof}

\begin{theorem}\label{thm:conj-unit-mul-one}
  \lean{KummerCriterion.conj_unit_mul_eq_one}\leanok
  \uses{thm:conj-generator-associated}
The unit $u$ obtained from $\overline a=ua$ in
\cref{thm:conj-generator-associated} satisfies
\[
  u\overline u=1.
\]
\end{theorem}
\begin{proof}\leanok
Apply complex conjugation to $\overline a=ua$: since conjugation is an
involution, $a = \overline u\,\overline a = \overline u u a$.  As $a\ne0$,
cancellation in the domain $\mathcal O_K$ gives $u\overline u=1$.
\end{proof}

\begin{theorem}\label{thm:antisymmetric-unit-classification}
  \lean{KummerCriterion.antisymmetric_unit_eq_neg_one_pow_mul_zeta_pow}\leanok
  \uses{thm:conj-unit-mul-one}
Let $K$ be the $p$th cyclotomic field with $p$ odd.  If
$u\in\mathcal O_K^\times$ satisfies $u\overline u=1$, then
\[
  u=(-1)^n\zeta_p^m
\]
for some integers $n,m$.
\end{theorem}
\begin{proof}\leanok
At every complex embedding $\sigma$ of $K$, conjugation corresponds to
complex conjugation of the image, so $|\sigma(u)|^2 = \sigma(u
\overline u) = 1$.  An algebraic integer all of whose conjugates have
absolute value $1$ is a root of unity (Kronecker's theorem), so $u$ is a
torsion unit.  The torsion subgroup of $\Qzetap^\times$ for $p$ odd is
generated by $-1$ and $\zeta_p$, which gives the displayed form.
\end{proof}

\begin{theorem}\label{thm:conj-fixed-associate}
  \lean{KummerCriterion.exists_conj_fixed_associate_of_classification}\leanok
  \uses{thm:antisymmetric-unit-classification}
In the situation of \cref{thm:conj-generator-associated}, the generator $a$
may be replaced by an associate $b$ such that
\[
  \overline b=b
  \qquad\text{and}\qquad
  (b)=(a).
\]
\end{theorem}
\begin{proof}\leanok
By \cref{thm:antisymmetric-unit-classification}, the unit $u =
\overline a / a$ has the form $(-1)^n \zeta_p^m$.  Since $p$ is odd, every
power of $\zeta_p$ is a square in $\langle\zeta_p\rangle$, and the sign can
be absorbed: write $u = w/\overline w$ for a suitable root of unity $w$
(concretely, a power of $\zeta_p$ chosen so that $w^2$ accounts for
$\zeta_p^m$, possibly times $-1$).  Then $b := w a$ satisfies
$\overline b = \overline w\,\overline a = \overline w u a = w a = b$,
and $b$ is an associate of $a$, so it generates the same ideal.
\end{proof}

\begin{theorem}\label{thm:mem-ring-of-integers-conj-self}
  \lean{KummerCriterion.mem_ringOfIntegers_of_conj_eq_self}\leanok
  \uses{def:cm-real-subfield}
If $b\in\mathcal O_K$ is fixed by complex conjugation, then it is the image
of an element of $\mathcal O_{K^+}$.
\end{theorem}
\begin{proof}\leanok
For a CM field, $K^+$ is exactly the fixed field of conjugation, so the
element $b$ lies in $K^+$.  Being integral over $\Z$, it lies in
$K^+ \cap \mathcal O_K = \mathcal O_{K^+}$.
\end{proof}

\begin{theorem}\label{thm:is-principal-after-extension}
  \lean{KummerCriterion.isPrincipal_of_isPrincipal_map_Kplus}\leanok
  \uses{thm:is-principal-map-span-singleton, thm:ideal-map-conj,
  thm:conj-generator-associated, thm:conj-unit-mul-one,
  thm:antisymmetric-unit-classification, thm:conj-fixed-associate,
  thm:mem-ring-of-integers-conj-self}
Let $K$ be the $p$th cyclotomic field, with $p$ odd.  If an ideal of
$\mathcal O_{K^+}$ becomes principal after extension to $\mathcal O_K$,
then it was already principal.
\end{theorem}
\begin{proof}\leanok
Let $I$ be such an ideal; we may assume $I \ne 0$.  Choose a generator $a$
of $I\mathcal O_K$.  The extended ideal is conjugation-stable
(\cref{thm:ideal-map-conj}), so \cref{thm:conj-generator-associated}
produces a unit $u$ with $\overline a = ua$, and
\cref{thm:conj-unit-mul-one} shows $u\overline u = 1$.  The classification
of such units (\cref{thm:antisymmetric-unit-classification}) lets us
replace $a$ by a conjugation-fixed associate $b$ with $(b) = I\mathcal O_K$
(\cref{thm:conj-fixed-associate}).  By
\cref{thm:mem-ring-of-integers-conj-self}, $b$ descends to an element $b_0
\in \mathcal O_{K^+}$, and \cref{thm:is-principal-map-span-singleton}
contracts the equality $I\mathcal O_K = (b_0)$ to $I = (b_0)$.
\end{proof}

\section{The relative class number}

\begin{definition}\label{def:class-group-map}
  \lean{KummerCriterion.classGroupMap}\leanok
  \uses{def:cm-real-subfield}
The inclusion $\mathcal O_{K^+}\subseteq\mathcal O_K$ induces, by extension
of fractional ideals, the class-group homomorphism
\[
  \Cl{\mathcal O_{K^+}}\longrightarrow \Cl{\mathcal O_K}.
\]
\end{definition}

\begin{theorem}\label{thm:class-group-map-injective}
  \lean{KummerCriterion.classGroupMap_injective}\leanok
  \uses{def:class-group-map, thm:is-principal-after-extension}
For an odd prime cyclotomic field, the map
\[
  \Cl{\mathcal O_{K^+}}\longrightarrow \Cl{\mathcal O_K}
\]
is injective.
\end{theorem}
\begin{proof}\leanok
A group homomorphism is injective when its kernel is trivial.  If an ideal
class of $\mathcal O_{K^+}$ maps to the trivial class, a representative
ideal becomes principal after extension to $\mathcal O_K$; by
\cref{thm:is-principal-after-extension} the representative was already
principal, so the class is trivial.
\end{proof}

\begin{theorem}\label{thm:hplus-divides-h}
  \lean{KummerCriterion.hPlus_dvd_h}\leanok
  \uses{def:class-number-h, def:class-number-hplus,
  thm:class-group-map-injective}
The plus class number divides the full class number:
\[
  h^+(K)\mid h(K).
\]
\end{theorem}
\begin{proof}\leanok
By \cref{thm:class-group-map-injective}, $\Cl{\mathcal O_{K^+}}$ embeds as
a subgroup of the finite group $\Cl{\mathcal O_K}$, and Lagrange's theorem
gives the divisibility of cardinalities.
\end{proof}

\begin{definition}\label{def:class-number-hminus}
  \lean{KummerCriterion.hMinus}\leanok
  \uses{thm:hplus-divides-h}
The relative class number is the natural-number quotient
\[
  h^-(K)=h(K)/h^+(K),
\]
which is well defined by \cref{thm:hplus-divides-h}.
\end{definition}

\begin{theorem}\label{thm:h-factorisation}
  \lean{KummerCriterion.h_eq_hPlus_mul_hMinus}\leanok
  \uses{def:class-number-hminus, thm:hplus-divides-h}
The class number factors as
\[
  h(K)=h^+(K)\,h^-(K).
\]
\end{theorem}
\begin{proof}\leanok
Multiply the defining quotient of \cref{def:class-number-hminus} by
$h^+(K)$; this is legitimate in $\N$ because $h^+(K) \mid h(K)$
(\cref{thm:hplus-divides-h}) and $h^+(K) \ne 0$.
\end{proof}

\begin{theorem}\label{thm:dvd-h-iff-dvd-hminus-from-plus-to-minus}
  \lean{KummerCriterion.dvd_h_iff_dvd_hMinus_of_dvd_hPlus_imp}\leanok
  \uses{thm:h-factorisation}
Assume that a separate argument proves the implication
\[
  p\mid h^+(K)\Longrightarrow p\mid h^-(K).
\]
Then
\[
  p\mid h(K)\Longleftrightarrow p\mid h^-(K).
\]
\end{theorem}
\begin{proof}\leanok
Write $h(K) = h^+(K)\,h^-(K)$ by \cref{thm:h-factorisation}.  If $p \mid
h(K)$, then by primality $p\mid h^+(K)$ or $p\mid h^-(K)$, and the first
alternative implies the second by hypothesis; either way $p \mid h^-(K)$.
Conversely $h^-(K)$ divides $h(K)$, so $p \mid h^-(K)$ implies
$p \mid h(K)$.
\end{proof}


\pagebreak
\chapter{The minus class-number criterion}\label{ch:hminus-formula}
\section[Analytic class number formula for h-]{Analytic class number formula for $\hminus$}

Let
\[
  K=\Q(\zeta_p),
  \qquad
  K^+=\Q(\zeta_p)^+,
  \qquad
  n=\frac{p-1}{2}.
\]
Then $K$ has signature $(0,n)$ and $K^+$ has signature $(n,0)$.

\subsection*{The odd $L$-product as a residue quotient}

The factorisation of the cyclotomic Dedekind zeta function into Dirichlet
$L$-functions gives
\[
  \zeta_K(s)=\prod_{\chi \bmod p} L(s,\chi),
\]
where the product runs over all characters modulo $p$. The even characters are
exactly the characters of \(\Gal(K^+/\Q)\), so similarly
\[
  \zeta_{K^+}(s)=\prod_{\substack{\chi\bmod p\\ \chi\text{ even}}} L(s,\chi).
\]
Dividing yields
\[
  \frac{\zeta_K(s)}{\zeta_{K^+}(s)}=\prod_{\chi\text{ odd}}L(s,\chi).
\]
Since every odd character modulo $p$ is non-trivial, every factor on the right
is holomorphic at $s=1$. Taking residues at $s=1$ therefore gives:

\begin{proposition}\label{prop:odd-L-product-residue}
  \lean{BernoulliRegular.residue_ready_factorization_even_odd}
  \leanok
  \uses{def:cm-real-subfield, def:dirichlet-character, def:even-odd}
One has
\[
  \prod_{\chi\text{ odd}} L(1,\chi)
  = \frac{\operatorname*{Res}_{s=1}\zeta_K(s)}{\operatorname*{Res}_{s=1}\zeta_{K^+}(s)}.
\]
\end{proposition}

\begin{proof}\leanok
The quotient \(\zeta_K(s)/\zeta_{K^+}(s)\) has a removable singularity at
$s=1$, because both numerator and denominator have simple poles there. Its
value at $s=1$ is therefore the ratio of the residues. On the other hand this
quotient is exactly the product of the odd Dirichlet $L$-functions, all of
which are regular at $s=1$.
\end{proof}

\subsection*{The relative class number formula}

The analytic class number formula gives
\[
  \operatorname*{Res}_{s=1}\zeta_F(s)
  = \frac{2^{r_1(F)}(2\pi)^{r_2(F)}h_F R_F}{w_F\sqrt{|D_F|}}
\]
for every number field $F$, where $R_F$ is the regulator, $w_F$ the number of
roots of unity, and $D_F$ the discriminant.

For the cyclotomic fields under consideration one uses the standard identities
\[
  w_K=2p,
  \qquad
  w_{K^+}=2,
  \qquad
  R_K = 2^{n-1}R_{K^+},
  \qquad
  |D_K|=p^n|D_{K^+}|.
\]
Combining these with the definition \(h^- = h_K/h_{K^+}\) gives:

\begin{proposition}\label{prop:relative-class-number-L-product}
  \lean{BernoulliRegular.hMinus_formula_via_residues}
  \lean{BernoulliRegular.hMinus_formula_of_residue_and_hPlus_cyclotomic_and_gauss}
  \leanok
  \uses{prop:odd-L-product-residue, thm:h-factorisation, def:class-number-h,
    def:class-number-hplus, def:class-number-hminus}
One has
\[
  h^- = \frac{2p\,p^{n/2}}{(2\pi)^n}\prod_{\chi\text{ odd}}L(1,\chi).
\]
\end{proposition}

\begin{proof}\leanok
Applying the analytic class number formula to $K$ and $K^+$ gives
\[
  h_K = \operatorname*{Res}_{s=1}\zeta_K(s)\cdot
  \frac{w_K\sqrt{|D_K|}}{(2\pi)^n R_K}
\]
and
\[
  h_{K^+} = \operatorname*{Res}_{s=1}\zeta_{K^+}(s)\cdot
  \frac{w_{K^+}\sqrt{|D_{K^+}|}}{2^n R_{K^+}}.
\]
Taking the quotient,
\[
  h^- =
  \frac{\operatorname*{Res}_{s=1}\zeta_K(s)}{\operatorname*{Res}_{s=1}\zeta_{K^+}(s)}
  \cdot \frac{w_K}{w_{K^+}}
  \cdot \frac{2^n R_{K^+}}{(2\pi)^n R_K}
  \cdot \sqrt{\frac{|D_K|}{|D_{K^+}|}}.
\]
Substituting the four identities above gives
\[
  h^- = \frac{2p\,p^{n/2}}{(2\pi)^n}
  \frac{\operatorname*{Res}_{s=1}\zeta_K(s)}{\operatorname*{Res}_{s=1}\zeta_{K^+}(s)}.
\]
Now apply \cref{prop:odd-L-product-residue}.
\end{proof}

\subsection*{The product of the odd Gauss sums}

For a primitive character $\chi$ modulo $p$, write
\[
  \tau(\chi)=\sum_{a\in \Z/p\Z}\chi(a)e^{2\pi i a/p}.
\]
The standard identity
\[
  \tau(\chi)\tau(\chi^{-1}) = \chi(-1)p
\]
shows that each inverse pair of odd characters contributes \(-p\) to the
product of all odd Gauss sums.

\begin{lemma}\label{lem:odd-gauss-product}
  \lean{BernoulliRegular.rawGaussProduct}
  \lean{BernoulliRegular.cyclotomicHGaussGoal_holds}
  \leanok
  \uses{def:gauss-sum, prop:gauss-norm, def:even-odd}
One has
\[
  \prod_{\chi\text{ odd}} \tau(\chi) = i^n p^{n/2}.
\]
\end{lemma}

\begin{proof}\leanok
If $p\equiv 1\pmod 4$, then the quadratic character is even, so no odd
character is self-inverse. The odd characters split into \(n/2\) inverse
pairs, each contributing \(-p\). Therefore
\[
  \prod_{\chi\text{ odd}}\tau(\chi)=(-p)^{n/2}=i^n p^{n/2}.
\]

If $p\equiv 3\pmod 4$, then the quadratic character \(\lambda\) is odd and is
the unique self-inverse odd character. The remaining odd characters split into
\((n-1)/2\) inverse pairs, so
\[
  \prod_{\chi\text{ odd}}\tau(\chi)
  = \tau(\lambda)(-p)^{(n-1)/2}.
\]
The classical quadratic Gauss sum evaluation gives
\(
  \tau(\lambda)=i\sqrt p.
\)
Hence
\[
  \prod_{\chi\text{ odd}}\tau(\chi)
  = i\sqrt p\,(-p)^{(n-1)/2}
  = i^n p^{n/2}.
\]
\end{proof}

\begin{theorem}[Analytic formula for \texorpdfstring{$\hminus$}{h-}]
  \label{thm:hminus-formula}
  \lean{BernoulliRegular.hMinus_formula}
  \leanok
  \uses{prop:relative-class-number-L-product, cor:L1-odd, lem:odd-gauss-product,
    def:gen-bernoulli}
The relative class number satisfies
\[
  h^- = 2p\prod_{\chi\text{ odd}}
  \left(-\frac12\,\Bgen{1}{\chi^{-1}}\right).
\]
\end{theorem}

\begin{proof}\leanok
There are exactly \(n=(p-1)/2\) odd characters modulo $p$. By
\cref{prop:relative-class-number-L-product},
\[
  h^- = \frac{2p\,p^{n/2}}{(2\pi)^n}
  \prod_{\chi\text{ odd}}L(1,\chi).
\]
For each odd character \(\chi\), \cref{cor:L1-odd} gives
\[
  L(1,\chi)=\frac{\pi i\,\tau(\chi)}{p}\,\Bgen{1}{\chi^{-1}}.
\]
Multiplying over all odd characters yields
\[
  \prod_{\chi\text{ odd}}L(1,\chi)
  = \left(\frac{\pi i}{p}\right)^n
    \left(\prod_{\chi\text{ odd}}\tau(\chi)\right)
    \left(\prod_{\chi\text{ odd}}\Bgen{1}{\chi^{-1}}\right).
\]
By \cref{lem:odd-gauss-product}, this becomes
\[
  \prod_{\chi\text{ odd}}L(1,\chi)
  = \frac{\pi^n i^{2n}}{p^{n/2}}
    \prod_{\chi\text{ odd}}\Bgen{1}{\chi^{-1}}
  = \frac{\pi^n(-1)^n}{p^{n/2}}
    \prod_{\chi\text{ odd}}\Bgen{1}{\chi^{-1}}.
\]
Substituting back gives
\[
  h^- = \frac{2p\,p^{n/2}}{(2\pi)^n}
        \cdot \frac{\pi^n(-1)^n}{p^{n/2}}
        \prod_{\chi\text{ odd}}\Bgen{1}{\chi^{-1}}
      = 2p\left(-\frac12\right)^n
        \prod_{\chi\text{ odd}}\Bgen{1}{\chi^{-1}}.
\]
Since there are exactly \(n\) odd characters, this is precisely
\[
  2p\prod_{\chi\text{ odd}}
  \left(-\frac12\,\Bgen{1}{\chi^{-1}}\right).
\]
\end{proof}


\pagebreak
\chapter{Cyclotomic units and Kummer logarithms}\label{ch:cyclotomic-units}
% Cyclotomic units and the plus class number.

This chapter records the cyclotomic-unit part of the proof.  The purpose is
to connect Bernoulli non-divisibility with the plus class number. Throughout
this chapter $p$ is an odd prime, $K$ is a field isomorphic to
$\Q(\zeta_p)$, and
\[
  K^+ = \Q(\zeta_p)^+
\]
is its maximal real subfield.  The group of real units is
\[
  E^+ := (\mathcal O_{K^+})^\times .
\]
Thus \(E^+\) is the full unit group of the maximal real subfield.

\section{Real cyclotomic units}

Let
\[
  u_a = \frac{1-\zeta_p^a}{1-\zeta_p}
  \qquad (a \in \Z,\ (a,p)=1)
\]
be the usual cyclotomic unit in $\mathcal O_K^\times$.  For the real
subfield one uses the conjugation-fixed unit
\[
  \varepsilon_a = u_a\,\overline{u_a}
  =
  \zeta_p^{\,1-a}
  \left(\frac{1-\zeta_p^a}{1-\zeta_p}\right)^2 .
\]
It lies in $K^+$ because complex conjugation fixes it, and it is a unit
because both $u_a$ and $\overline{u_a}$ are units.  The range used for prime
conductor is
\[
  2 \le a \le \frac{p-1}{2}.
\]
In that range $0<a<p$, so a prime $p$ cannot divide $a$; hence $(a,p)=1$ and
the above cyclotomic unit is defined.

\begin{lemma}[Coprimality of the standard indices]
  \label{lem:real-cyclotomic-index-coprime}
  \lean{BernoulliRegular.realCyclotomicUnit_index_coprime}
  \leanok
  \uses{def:galois-cyclotomic}
Let $2 \le a \le (p-1)/2$. Then $a$ is coprime to $p$.
\end{lemma}

\begin{proof}\leanok
The lower bound gives $a>0$.  The upper bound gives
$a \le (p-1)/2 < p$, so $a<p$.  If $p$ divided $a$, then since $a$ is
positive we would have $p \le a$, contradicting $a<p$.  Because $p$ is prime,
not being divisible by $p$ is equivalent to being coprime to $p$.
\end{proof}

\begin{definition}[Real cyclotomic unit]
  \label{def:real-cyclotomic-unit}
  \lean{BernoulliRegular.realCyclotomicUnit}
  \leanok
  \uses{def:cm-real-subfield, lem:real-cyclotomic-index-coprime}
For $2 \le a \le (p-1)/2$, the unit
\[
  \varepsilon_a \in E^+
\]
is the unique real unit whose image in $\mathcal O_K^\times$ is the
conjugation-fixed product $u_a\overline{u_a}$.
\end{definition}

After extension of scalars from \(\mathcal O_{K^+}\) to \(\mathcal O_K\),
this unit is exactly the standard conjugation-fixed product
\(u_a\overline{u_a}\), and is fixed by complex conjugation.

\section{The cyclotomic-unit subgroups}

There are two closely related real cyclotomic-unit subgroups.  The logarithmic
saturation argument is most convenient for the squared
family $\varepsilon_a$.  The index theorem is naturally stated for the
normalised family
\[
  \eta_a
  =
  \zeta_p^{(1-a)/2}\frac{1-\zeta_p^a}{1-\zeta_p},
  \qquad
  \eta_a^2=\varepsilon_a .
\]
Here $(1-a)/2$ is understood modulo $p$, using that $2$ is invertible modulo
the odd prime $p$.

\begin{definition}[The squared cyclotomic-unit subgroup]
  \label{def:CPlus-squared}
  \lean{BernoulliRegular.CPlus}
  \leanok
  \uses{def:real-cyclotomic-unit}
Let $g=(p-3)/2$.  The subgroup $C^+_{\mathrm{sq}}\le E^+$ is generated by
$-1$ and by the units
\[
  \varepsilon_2,\varepsilon_3,\ldots,\varepsilon_{(p-1)/2}.
\]
Equivalently,
\[
  C^+_{\mathrm{sq}}
  =
  \left\langle
    -1,\ \varepsilon_a\ (2\le a\le (p-1)/2)
  \right\rangle .
\]
\end{definition}

The finite index set may be written as
\[
  i\in \operatorname{Fin}((p-3)/2), \qquad a=i+2.
\]
Thus the listed units are precisely the nontrivial real cyclotomic generators
in the prime-conductor range, together with the torsion unit \(-1\).

\begin{definition}[The normalised cyclotomic-unit subgroup]
  \label{def:CPlus-normalized}
  \lean{BernoulliRegular.normalizedCPlus}
  \leanok
  \uses{def:real-cyclotomic-unit}
The subgroup $C^+_{\mathrm{norm}}\le E^+$ is generated by $-1$ and by the
normalised real cyclotomic units
\[
  \eta_2,\eta_3,\ldots,\eta_{(p-1)/2}.
\]
\end{definition}

The relation between the two families is exactly the relation
$\eta_a^2=\varepsilon_a$.

\begin{lemma}[Squares of normalised generators]
  \label{lem:normalized-square}
  \lean{BernoulliRegular.normalizedCPlusGenerator_sq_eq_CPlusGenerator}
  \leanok
  \uses{def:CPlus-squared, def:CPlus-normalized}
For each standard generator index $a$,
\[
  \eta_a^2=\varepsilon_a.
\]
\end{lemma}

\begin{proof}\leanok
This is the elementary identity
\[
  \left(
    \zeta_p^{(1-a)/2}\frac{1-\zeta_p^a}{1-\zeta_p}
  \right)^2
  =
  \zeta_p^{1-a}
  \left(\frac{1-\zeta_p^a}{1-\zeta_p}\right)^2
  =
  u_a\overline{u_a}.
\]
The last equality is the standard computation
\[
  \overline{u_a}
  =
  \frac{1-\zeta_p^{-a}}{1-\zeta_p^{-1}}
  =
  \zeta_p^{1-a}\frac{1-\zeta_p^a}{1-\zeta_p}.
\]
\end{proof}

\begin{lemma}[Odd-primary comparison of the two indices]
  \label{lem:CPlus-index-normalized-index}
  \lean{BernoulliRegular.CPlus_index_prime_dvd_iff_normalizedCPlus_index_prime_dvd}
  \lean{BernoulliRegular.subgroup_index_prime_dvd_iff_of_relIndex_dvd_two_pow}
  \leanok
  \uses{def:CPlus-squared, def:CPlus-normalized, lem:normalized-square}
For odd $p$,
\[
  p \mid [E^+ : C^+_{\mathrm{sq}}]
  \quad\Longleftrightarrow\quad
  p \mid [E^+ : C^+_{\mathrm{norm}}].
\]
\end{lemma}

\begin{proof}\leanok
The subgroup $C^+_{\mathrm{sq}}$ is contained in
$C^+_{\mathrm{norm}}$, because $-1$ is a normalised generator and
each squared generator $\varepsilon_a$ is the square of the normalised
generator $\eta_a$.  Conversely, every class in
$C^+_{\mathrm{norm}}/C^+_{\mathrm{sq}}$ has square one: it is enough to check
this on the generators, and it follows from $\eta_a^2=\varepsilon_a$ and
$(-1)^2=1$.  Hence the relative index
$[C^+_{\mathrm{norm}}:C^+_{\mathrm{sq}}]$ divides a power of $2$.

The two ambient indices differ by this relative index.  Since $p$ is odd, $p$
does not divide any power of $2$.  Therefore the $p$-primary divisibility of
$[E^+ : C^+_{\mathrm{sq}}]$ and $[E^+ : C^+_{\mathrm{norm}}]$ is the same.
\end{proof}

\section{Exact p-saturation}

Let $G$ be an abelian group and $H\le E\le G$.  We use the exact image of the
$p$-th power map:
\[
  H^p := \{h^p : h\in H\}.
\]
This is not a closure or saturation of the image; it is the literal subgroup
of $p$-th powers.

\begin{definition}[Exact $p$-saturation]
  \label{def:p-saturated}
  \lean{BernoulliRegular.pSaturated}
  \lean{BernoulliRegular.pPowerSubgroup}
  \leanok
The subgroup $H$ is $p$-saturated in $E$ if $H\le E$ and
\[
  H\cap E^p \subseteq H^p.
\]
Equivalently, whenever $x\in H$ and $x=y^p$ for some $y\in E$, there exists
$z\in H$ such that $x=z^p$.
\end{definition}

The useful criterion for $C^+_{\mathrm{sq}}$ is stated in terms of exponent
vectors.  Every element of $C^+_{\mathrm{sq}}$ can be written as
\[
  (-1)^s\prod_{a=2}^{(p-1)/2}\varepsilon_a^{e_a},
  \qquad s,e_a\in\Z .
\]
The following criterion uses these exponent vectors.

\begin{theorem}[Exponent criterion for saturation]
  \label{thm:CPlus-pSaturated-exponents}
  \lean{BernoulliRegular.exists_CPlusExponentProduct_of_mem_CPlus}
  \lean{BernoulliRegular.CPlus_pSaturated_of_generator_exponents_modP_zero}
  \lean{BernoulliRegular.CPlusExponentProduct_mem_CPlus}
  \lean{BernoulliRegular.neg_one_zpow_mem_pPowerSubgroup_CPlus}
  \lean{BernoulliRegular.CPlusExponentProduct_pow_of_exponents_eq_mul}
  \leanok
  \uses{def:p-saturated, def:CPlus-squared}
Assume that for every expression
\[
  x=(-1)^s\prod_a \varepsilon_a^{e_a}\in C^+_{\mathrm{sq}}
\]
which is a $p$-th power in $E^+$, all exponents satisfy
\[
  e_a\equiv 0 \pmod p.
\]
Then $C^+_{\mathrm{sq}}$ is $p$-saturated in $E^+$.
\end{theorem}

\begin{proof}\leanok
Let $x\in C^+_{\mathrm{sq}}$ and suppose $x=y^p$ for some $y\in E^+$.
Choose an exponent expression
\[
  x=(-1)^s\prod_a \varepsilon_a^{e_a}.
\]
By hypothesis, $e_a\equiv 0\pmod p$ for every $a$, so choose integers
$k_a$ with $e_a=pk_a$.  Since $p$ is odd, the sign term is already a
$p$-th power of itself:
\[
  ((-1)^s)^p=(-1)^s.
\]
Therefore
\[
  x
  =
  (-1)^s\prod_a \varepsilon_a^{pk_a}
  =
  \left((-1)^s\prod_a \varepsilon_a^{k_a}\right)^p .
\]
The element inside the parentheses belongs to $C^+_{\mathrm{sq}}$, so $x$ is
a $p$-th power inside $C^+_{\mathrm{sq}}$.  This is exactly
$p$-saturation.
\end{proof}

\section{From saturation to index non-divisibility}

\begin{theorem}[Saturation forbids $p$ in the index]
  \label{thm:not-dvd-index-of-pSaturated}
  \lean{BernoulliRegular.not_dvd_index_of_pSaturated}
  \lean{BernoulliRegular.subgroup_not_dvd_index_of_pSaturated_top_of_pow_eq_one_mem}
  \lean{BernoulliRegular.CPlus_index_ne_zero}
  \leanok
  \uses{def:p-saturated, def:CPlus-squared}
If $C^+_{\mathrm{sq}}$ is $p$-saturated in $E^+$, then
\[
  p\nmid [E^+:C^+_{\mathrm{sq}}].
\]
\end{theorem}

\begin{proof}\leanok
The proof is purely group-theoretic once one knows that
$C^+_{\mathrm{sq}}$ has finite index in $E^+$ and contains the torsion of
$E^+$.  Suppose for contradiction that $p$ divides
$[E^+:C^+_{\mathrm{sq}}]$.  By Cauchy's theorem applied to the finite quotient
$E^+/C^+_{\mathrm{sq}}$, there is a non-trivial quotient class of order $p$.
Choose a representative $g\in E^+$.  The order condition says
\[
  g^p\in C^+_{\mathrm{sq}}.
\]
Of course $g^p$ is a $p$-th power in $E^+$.  By $p$-saturation, there is
$h\in C^+_{\mathrm{sq}}$ with
\[
  h^p=g^p.
\]
Then
\[
  (gh^{-1})^p=1.
\]
Thus $gh^{-1}$ is $p$-torsion in $E^+$.  The only torsion units in the real
cyclotomic field are $\pm1$, and both belong to $C^+_{\mathrm{sq}}$.  Hence
$gh^{-1}\in C^+_{\mathrm{sq}}$, and since $h\in C^+_{\mathrm{sq}}$, also
$g\in C^+_{\mathrm{sq}}$.  This means the quotient class of $g$ was trivial,
contradicting its order $p$.
\end{proof}

\section{The Kummer logarithm determinant}

The $p$-adic part of the argument studies the logarithms of the real
cyclotomic units at the prime above $p$.  After passing to the completed local
ring and reducing suitable logarithmic coordinates modulo $p$, one obtains a
matrix over $\mathbb F_p$.  Its columns are indexed by the generators
\[
  \varepsilon_2,\ldots,\varepsilon_{(p-1)/2},
\]
and its rows are indexed by the even powers which occur in the Kummer range.

The determinant has two parts.  One part is a finite-field Vandermonde
determinant in the Teichmuller nodes; this is nonzero because the nodes are
distinct.  The other part is the product of Bernoulli factors.  Consequently
the determinant is nonzero exactly when no Bernoulli numerator in Kummer's
range is divisible by $p$.

\begin{theorem}[Kummer logarithm determinant criterion]
  \label{thm:kummer-log-det-bernoulli}
  \lean{BernoulliRegular.CyclotomicUnits.kummerLogMatrix_det_ne_zero_iff_bernoulli_nonzero}
  \lean{BernoulliRegular.CyclotomicUnits.concreteKummerLogMatrix_eq_diagonal_mul_vandermonde}
  \lean{BernoulliRegular.CyclotomicUnits.concreteKummerLogMatrix_det_ne_zero_iff_forall_bernoulliFactor_ne_zero}
  \lean{BernoulliRegular.CyclotomicUnits.bernoulliFactor_ne_zero_iff_not_dvd_bernoulli_num}
  \lean{BernoulliRegular.CyclotomicUnits.vandermonde_teichmuller_even_sub_one_det_ne_zero}
  \leanok
  \uses{def:CPlus-squared, thm:bernoulli-padicInt-below, cor:bernoulli-den-coprime}
Assume $p\ge 5$.  The Kummer logarithm matrix has nonzero determinant if and
only if
\[
  \forall j,\quad
  1\le j,\ 2j\le p-3
  \Longrightarrow
  p\nmid (B_{2j})_{\mathrm{num}}.
\]
\end{theorem}

\begin{proof}\leanok
The determinant computation factors the matrix as
\[
  \operatorname{diag}(r_1,\ldots,r_g)\,V,
  \qquad g=(p-3)/2,
\]
where $V$ is the finite-field Vandermonde matrix attached to the even
Teichmuller nodes and $r_j$ is the Bernoulli row factor attached to
$B_{2j}/(2j)$.  The determinant is therefore
\[
  \det(V)\prod_j r_j.
\]
The Vandermonde determinant is nonzero because the Teichmuller nodes are
pairwise distinct in $\mathbb F_p$.  Thus the determinant is nonzero exactly
when every $r_j$ is nonzero.

For $1\le j$ and $2j\le p-3$, the denominator of $B_{2j}$ is prime to $p$ by
the integrality results for Bernoulli numbers below the boundary.  Also
$2j<p$.  Hence reducing $B_{2j}/(2j)$ modulo $p$ gives zero exactly when the
integer numerator of $B_{2j}$ is divisible by $p$.  This identifies the
non-vanishing of all row factors with the displayed Bernoulli
non-divisibility condition.
\end{proof}

\begin{theorem}[Determinant non-vanishing gives saturation]
  \label{thm:kummer-log-det-saturation}
  \lean{BernoulliRegular.CyclotomicUnits.cyclotomicUnits_pSaturated_of_kummerLog_det_ne_zero}
  \lean{BernoulliRegular.CyclotomicUnits.completedLog_relation_of_CPlus_product_mem_powers}
  \lean{BernoulliRegular.CyclotomicUnits.concreteKummerLogMatrix_mulVec_exponents_eq_zero}
  \lean{BernoulliRegular.CyclotomicUnits.exponents_modP_eq_zero_of_kummerLogMatrix_relation}
  \lean{BernoulliRegular.CyclotomicUnits.CPlusGenerator_exponents_modP_zero_of_kummerLog_det_ne_zero}
  \leanok
  \uses{def:p-saturated, def:CPlus-squared, thm:CPlus-pSaturated-exponents}
Assume $p\ge 5$.  If the Kummer logarithm determinant is nonzero, then
$C^+_{\mathrm{sq}}$ is $p$-saturated in $E^+$.
\end{theorem}

\begin{proof}\leanok
Suppose an element
\[
  x=(-1)^s\prod_a \varepsilon_a^{e_a}
\]
is a $p$-th power in $E^+$.  Taking completed $p$-adic logarithms of the
local image of this relation gives an additive relation
\[
  \sum_a e_a\,\log(\varepsilon_a)
  \in p\cdot L,
\]
where $L$ is the completed logarithm lattice. The
sign term disappears from the logarithmic coordinates, because $p$ is odd and
the sign is torsion.

Reducing the coordinates of this relation modulo $p$ gives
\[
  M\,(e_a\bmod p)_a=0
\]
for the Kummer logarithm matrix $M$.  If $\det(M)\ne0$, then $M$ is invertible
over $\mathbb F_p$, so every coordinate $e_a\bmod p$ is zero.  The exponent
criterion of \cref{thm:CPlus-pSaturated-exponents} then proves
$p$-saturation.
\end{proof}

\section{Bernoulli non-divisibility and the plus index}

The previous sections combine into the theorem used by the final proof.

\begin{theorem}[Bernoulli non-divisibility implies index non-divisibility]
  \label{thm:bernoulli-nonzero-index}
  \lean{BernoulliRegular.not_dvd_cyclotomicUnitIndex_of_bernoulli_nonzero}
  \leanok
  \uses{thm:kummer-log-det-bernoulli, thm:kummer-log-det-saturation,
    thm:not-dvd-index-of-pSaturated}
Assume $p$ is odd and
\[
  \forall j,\quad
  1\le j,\ 2j\le p-3
  \Longrightarrow
  p\nmid (B_{2j})_{\mathrm{num}}.
\]
Then
\[
  p\nmid [E^+:C^+_{\mathrm{sq}}].
\]
\end{theorem}

\begin{proof}\leanok
If $p=3$, then the generator range is empty.  The real subfield of
$\Q(\zeta_3)$ is $\Q$, so its unit group is generated by $-1$.  Hence
$C^+_{\mathrm{sq}}=E^+$ and the index is $1$, which is not divisible by $p$.

Now assume $p\ge5$.  The Bernoulli non-divisibility hypothesis makes the
Kummer logarithm determinant nonzero by
\cref{thm:kummer-log-det-bernoulli}.  The determinant non-vanishing theorem
\cref{thm:kummer-log-det-saturation} gives $p$-saturation of
$C^+_{\mathrm{sq}}$ in $E^+$.  Finally
\cref{thm:not-dvd-index-of-pSaturated} converts this saturation into
\[
  p\nmid [E^+:C^+_{\mathrm{sq}}].
\]
\end{proof}

\section{The needed Sinnott index formula}

The index theorem used in the final argument is the odd-primary consequence
of Sinnott's formula for real cyclotomic units.  The family used in the
regulator computation is the same squared family as above, but with the
torsion subgroup adjoined explicitly.

\begin{definition}[Sinnott's full-rank subgroup]
  \label{def:sinnott-index-subgroup}
  \lean{BernoulliRegular.cyclotomicUnitIndexSubgroup}
  \leanok
  \uses{def:real-cyclotomic-unit, def:CPlus-squared}
For $p\ge3$, let
\[
  C^+_{\mathrm{Sin}}
  =
  \left\langle
    \varepsilon_2,\varepsilon_3,\ldots,\varepsilon_{(p-1)/2}
  \right\rangle \vee E^+_{\mathrm{tors}}
  \le E^+ .
\]
Equivalently, $C^+_{\mathrm{Sin}}$ is generated by the full-rank real
cyclotomic-unit family used in the regulator determinant, together with all
torsion units of $E^+$.
\end{definition}

\begin{lemma}[Agreement with the squared subgroup]
  \label{lem:sinnott-subgroup-eq-CPlus}
  \lean{BernoulliRegular.cyclotomicUnitIndexSubgroup_eq_CPlus}
  \lean{BernoulliRegular.CPlusGenerator_eq_cyclotomicUnitFamilyKplus}
  \lean{BernoulliRegular.range_cyclotomicUnitFamilyKplusFinRank_eq}
  \lean{BernoulliRegular.torsionKplus_le_CPlus}
  \lean{BernoulliRegular.CPlus_le_cyclotomicUnitIndexSubgroup}
  \lean{BernoulliRegular.cyclotomicUnitIndexSubgroup_le_CPlus}
  \leanok
  \uses{def:sinnott-index-subgroup, def:CPlus-squared}
For $p\ge3$,
\[
  C^+_{\mathrm{Sin}}=C^+_{\mathrm{sq}}.
\]
\end{lemma}

\begin{proof}\leanok
The full-rank family in $C^+_{\mathrm{Sin}}$ is indexed by
$\operatorname{Fin}((p-3)/2)$, and the element with index $i$ is exactly
$\varepsilon_{i+2}$.  Thus its range is the same generator range used for
$C^+_{\mathrm{sq}}$.

It remains to compare torsion.  The torsion of the real cyclotomic unit group
is $\{\pm1\}$.  The unit $1$ belongs to every subgroup, and $-1$ is one of the
generators of $C^+_{\mathrm{sq}}$.  Hence adjoining torsion to the full-rank
family adds nothing beyond the subgroup already generated by
$-1,\varepsilon_2,\ldots,\varepsilon_{(p-1)/2}$.
\end{proof}

\begin{theorem}[Sinnott formula, odd-primary consequence]
  \label{thm:sinnott-index-pprimary}
  \lean{BernoulliRegular.cyclotomicUnitIndex_primeConductor_pPrimary_of_sinnottIndexFormula}
  \lean{BernoulliRegular.cyclotomicUnitIndex_primeConductor_pPrimary_of_kummerDirichletDeterminant}
  \lean{BernoulliRegular.FLT37.Sinnott.sinnottRegulatorIdentity_iff_kummerDirichletDeterminant}
  \lean{BernoulliRegular.FLT37.Sinnott.sinnottIndexFormula_of_regulatorIdentity}
  \lean{BernoulliRegular.FLT37.Sinnott.index_eq_twoPow_mul_hPlus_of_sinnottIndexFormula}
  \lean{BernoulliRegular.prime_dvd_two_pow_mul_iff_dvd}
  \leanok
  \uses{def:sinnott-index-subgroup}
Assume the Kummer-Dirichlet determinant identity for the real
cyclotomic-unit logarithm matrix.  Then
\[
  p\mid [E^+:C^+_{\mathrm{Sin}}]
  \quad\Longleftrightarrow\quad
  p\mid h^+ .
\]
\end{theorem}

\begin{proof}\leanok
Write $g=(p-3)/2$ and let
\[
  R^+ = \operatorname{Reg}(E^+)
\]
be the regulator of the maximal real subfield.  The determinant identity is
the regulator identity
\[
  \operatorname{Reg}(\varepsilon_2,\ldots,\varepsilon_{(p-1)/2})
  =
  2^g h^+ R^+ .
\]

The passage from regulators to indices is the standard covolume comparison
for the logarithmic unit lattice.  Modulo torsion, the units $E^+$ form a
lattice of rank $g$, and the family
$\varepsilon_2,\ldots,\varepsilon_{(p-1)/2}$ spans a full-rank sublattice.
The covolume of the latter is its index in the former times the covolume of
the former:
\[
  \frac{
    \operatorname{Reg}(\varepsilon_2,\ldots,\varepsilon_{(p-1)/2})
  }{R^+}
  =
  [E^+:C^+_{\mathrm{Sin}}].
\]
Since $R^+>0$, the displayed regulator identity therefore gives Sinnott's
index formula in this normalization:
\[
  [E^+:C^+_{\mathrm{Sin}}]=2^g h^+ .
\]

Finally, $p$ is odd, so $p\nmid 2^g$.  Euclid's lemma gives
\[
  p\mid 2^g h^+
  \quad\Longleftrightarrow\quad
  p\mid h^+,
\]
and combining this with the index formula proves the stated equivalence.
\end{proof}

\begin{theorem}[Deleted-Fourier determinant identity]
  \label{thm:deleted-fourier-kummer-dirichlet}
  \lean{BernoulliRegular.kummerDirichletDeterminant_of_deletedFourier}
  \leanok
  \uses{def:sinnott-index-subgroup, def:dirichlet-character,
    def:real-cyclotomic-unit}
Assume $p\ge5$.  The logarithm matrix of the full-rank real cyclotomic-unit
family satisfies the Kummer-Dirichlet determinant identity.
\end{theorem}

\begin{proof}\leanok
The rows of the matrix are indexed by the non-trivial even characters modulo
$p$, and the columns by
\[
  \varepsilon_2,\ldots,\varepsilon_{(p-1)/2}.
\]
After the standard Fourier change of basis, the determinant separates into a
deleted Fourier determinant and the product of the even character factors.

The deleted Fourier determinant is nonzero: it is the determinant of the
character table with the trivial row and the redundant column removed, and
orthogonality of characters gives its inverse after multiplying by a power of
$p$.  The remaining product is the even part of the cyclotomic class-number
formula.  With the present squared-unit normalization, the regulator identity
therefore has the form
\[
  \operatorname{Reg}(\varepsilon_2,\ldots,\varepsilon_{(p-1)/2})
  =
  2^g h^+ R^+ ,
  \qquad g=(p-3)/2.
\]
This is precisely the Kummer-Dirichlet determinant identity used in the
index calculation.
\end{proof}

\begin{corollary}[Sinnott formula for prime conductor]
  \label{cor:sinnott-index-prime-conductor}
  \lean{BernoulliRegular.kummerDirichletDeterminant_of_deletedFourier}
  \lean{BernoulliRegular.cyclotomicUnitIndex_primeConductor_pPrimary_aux}
  \lean{BernoulliRegular.cyclotomicUnitIndex_primeConductor_pPrimary_of_kummerDirichletDeterminant}
  \leanok
  \uses{thm:sinnott-index-pprimary, thm:deleted-fourier-kummer-dirichlet}
If $p\ge5$, then
\[
  p\mid [E^+:C^+_{\mathrm{Sin}}]
  \quad\Longleftrightarrow\quad
  p\mid h^+ .
\]
\end{corollary}

\begin{proof}\leanok
For $p\ge5$, the deleted-Fourier determinant computation evaluates the
logarithm determinant of the real cyclotomic units.  The matrix rows are the
even non-trivial characters, the columns are the units
$\varepsilon_2,\ldots,\varepsilon_{(p-1)/2}$, and the determinant splits into
a nonzero Fourier determinant times the product of the character factors.
The class-number formula for the real subfield identifies this product with
$h^+R^+$, with the extra factor $2^g$ coming from using the squared
cyclotomic-unit family.  Thus the Kummer-Dirichlet determinant identity
holds, and \cref{thm:sinnott-index-pprimary} applies.
\end{proof}

\begin{theorem}[Prime-conductor cyclotomic-unit index theorem]
  \label{thm:cyclotomic-unit-index-hplus}
  \lean{BernoulliRegular.cyclotomicUnitIndex_primeConductor_pPrimary}
  \lean{BernoulliRegular.cyclotomicUnitIndex_primeConductor_pPrimary_of_eq_three}
  \lean{BernoulliRegular.cyclotomicUnitIndex_primeConductor_pPrimary_of_five_le}
  \leanok
  \uses{cor:sinnott-index-prime-conductor, lem:sinnott-subgroup-eq-CPlus,
    lem:CPlus-index-normalized-index}
For odd $p$,
\[
  p \mid [E^+:C^+_{\mathrm{norm}}]
  \quad\Longleftrightarrow\quad
  p \mid h^+.
\]
\end{theorem}

\begin{proof}\leanok
There are two cases.

If \(p=3\), then \(K^+=\Q\).  The plus class number is therefore \(h^+=1\).
The real unit rank is zero, so the normalised cyclotomic-unit subgroup is the
whole unit group.  Hence both sides are the assertion that \(3\) divides \(1\),
and both are false.

Now assume $p\ne3$.  Since $p$ is an odd prime, this gives $p\ge5$.  By
\cref{cor:sinnott-index-prime-conductor},
\[
  p\mid [E^+:C^+_{\mathrm{Sin}}]
  \quad\Longleftrightarrow\quad
  p\mid h^+ .
\]
The subgroup identification
\cref{lem:sinnott-subgroup-eq-CPlus} replaces $C^+_{\mathrm{Sin}}$ by
$C^+_{\mathrm{sq}}$.  Finally, the odd-primary comparison
\cref{lem:CPlus-index-normalized-index} replaces
$C^+_{\mathrm{sq}}$ by $C^+_{\mathrm{norm}}$, because the quotient between
the two generated subgroups has exponent dividing a power of $2$.  Since
$p$ is odd, this quotient does not affect $p$-divisibility of the ambient
index.  The displayed equivalence follows.
\end{proof}


\pagebreak
\chapter{Final assembly}\label{ch:kummer-final}
% Final assembly of Kummer's criterion.

This chapter combines the minus criterion of
Chapter~\ref{ch:hminus-formula} with the cyclotomic-unit results of
Chapter~\ref{ch:cyclotomic-units}.  The bridge is the implication
$p \mid h^+ \Rightarrow p \mid h^-$: if $p \nmid h^-$, the minus criterion
clears every Bernoulli numerator in Kummer's range, the Kummer logarithm
determinant is then nonzero, $C^+$ is $p$-saturated, and the
prime-conductor index theorem forces $p \nmid h^+$.

\section{\texorpdfstring{The implication from $h^+$ to $h^-$}{The implication from h-plus to h-minus}}

\begin{theorem}\label{thm:bernoulli-nonzero-index}
  \lean{KummerCriterion.not_dvd_cyclotomicUnitIndex_of_bernoulli_nonzero}\leanok
  \uses{thm:kummer-log-det-bernoulli, thm:kummer-log-det-saturation,
  thm:not-dvd-index-of-pSaturated}
Let $p$ be an odd prime.  Suppose that
\[
  \forall j,\ 1\le j,\ 2j\le p-3,\quad
  p\nmid\operatorname{num}(B_{2j}).
\]
Then
\[
  p\nmid[(\mathcal O_{K^+})^\times:C^+].
\]
\end{theorem}
\begin{proof}\leanok
If $p=3$, Kummer's range is empty, $K^+ = \Q$, the group $C^+$ is the full
unit group $\{\pm1\}$, and the index is $1$.  For $p\ge5$,
\cref{thm:kummer-log-det-bernoulli} identifies the Bernoulli nonvanishing
hypothesis with the nonvanishing of the Kummer logarithm determinant.  By
\cref{thm:kummer-log-det-saturation}, the subgroup $C^+$ is then
$p$-saturated in the full unit group, and
\cref{thm:not-dvd-index-of-pSaturated} concludes that the index is prime
to $p$.
\end{proof}

\begin{theorem}\label{thm:not-dvd-hplus-of-not-dvd-hminus-units}
  \lean{KummerCriterion.not_dvd_hPlus_of_not_dvd_hMinus_units}\leanok
  \uses{thm:hminus-nondiv-to-bernoulli-nondiv, thm:bernoulli-nonzero-index,
  thm:sinnott-index-pprimary, thm:CPlus-normalized-index}
If $p\nmid h^-(K)$, then $p\nmid h^+(K)$.
\end{theorem}
\begin{proof}\leanok
By \cref{thm:hminus-nondiv-to-bernoulli-nondiv}, the hypothesis
$p\nmid h^-(K)$ implies that no Bernoulli numerator in Kummer's range is
divisible by $p$, so \cref{thm:bernoulli-nonzero-index} gives
\[
  p\nmid[(\mathcal O_{K^+})^\times:C^+].
\]
Suppose for contradiction that $p\mid h^+(K)$.  The prime-conductor index
theorem (\cref{thm:sinnott-index-pprimary}) then gives $p \mid
[(\mathcal O_{K^+})^\times:C^+_{\mathrm{norm}}]$, and by
\cref{thm:CPlus-normalized-index} this is equivalent to $p$ dividing the
ordinary $C^+$ index --- a contradiction.
\end{proof}

\begin{theorem}\label{thm:hplus-to-hminus-units}
  \lean{KummerCriterion.weakReflection_dvd_hMinus_of_dvd_hPlus_units}\leanok
  \uses{thm:not-dvd-hplus-of-not-dvd-hminus-units}
If $p\mid h^+(K)$, then $p\mid h^-(K)$.
\end{theorem}
\begin{proof}\leanok
Contrapositive of \cref{thm:not-dvd-hplus-of-not-dvd-hminus-units}: if
$p\nmid h^-(K)$, that theorem gives $p\nmid h^+(K)$.
\end{proof}

\section{The total class-number criterion}

\begin{theorem}\label{thm:dvd-h-iff-bernoulli-units}
  \lean{KummerCriterion.dvd_h_iff_exists_dvd_bernoulli_units}\leanok
  \uses{thm:dvd-h-iff-dvd-hminus-from-plus-to-minus,
  thm:hplus-to-hminus-units, thm:hminus-bernoulli-final}
For an odd prime $p$ and the $p$th cyclotomic field $K$,
\[
  p\mid h(K)
  \quad\Longleftrightarrow\quad
  \exists k,\ 1\le k,\ 2k\le p-3,\quad
  p\mid\operatorname{num}(B_{2k}).
\]
\end{theorem}
\begin{proof}\leanok
The implication $p\mid h^+(K)\Rightarrow p\mid h^-(K)$ is
\cref{thm:hplus-to-hminus-units}; feeding it to
\cref{thm:dvd-h-iff-dvd-hminus-from-plus-to-minus} converts $p\mid h(K)$
into $p\mid h^-(K)$.  The minus class-number criterion
(\cref{thm:hminus-bernoulli-final}) rewrites the latter as the existence of
a Bernoulli numerator in Kummer's range divisible by $p$.
\end{proof}

\section{Kummer's criterion}

\begin{theorem}\label{thm:kummer-criterion}
  \lean{KummerCriterion}\leanok
  \uses{def:regular-prime, thm:dvd-h-iff-bernoulli-units}
Let $p$ be an odd prime.  Then
\[
  p\ \text{is regular}
  \quad\Longleftrightarrow\quad
  \forall k,\ 1\le k,\ 2k\le p-3,\quad
  p\nmid\operatorname{num}(B_{2k}).
\]
\end{theorem}
\begin{proof}\leanok
By \cref{def:regular-prime}, $p$ is regular if and only if $p$ is coprime
to the class number of the $p$th cyclotomic field; since $p$ is prime,
this is equivalent to $p\nmid h(K)$.  By
\cref{thm:dvd-h-iff-bernoulli-units}, the negation of $p\mid h(K)$ is
exactly the assertion that no $k$ in the range $1\le k$, $2k\le p-3$
satisfies $p\mid\operatorname{num}(B_{2k})$, which is the displayed
universal nondivisibility.
\end{proof}


\pagebreak
\chapter{Infinitely many irregular primes}\label{ch:irregular-primes}
% Carlitz's proof that there are infinitely many irregular primes.

This chapter records the Carlitz route from Kummer congruences to the
infinitude of irregular primes.  The theorem proved in Lean is
\[
  \{p : \N \mid p\text{ is prime and }p\text{ is not regular}\}
  \quad\text{is infinite}.
\]
The proof is by finite-set escape: given a finite set \(S\) of candidate
irregular primes, construct a new prime \(p\notin S\) dividing the numerator
of a suitable divided Bernoulli number \(B_M/M\), then use Kummer congruences
and Kummer's criterion to show that \(p\) is irregular.

\section{The divided Kummer congruence}

The Carlitz argument needs a Kummer congruence for the divided Bernoulli
numbers \(B_m/m\) with no upper bound on the exponents.

\begin{theorem}[Full divided Kummer congruence]
  \label{thm:full-divided-kummer}
  \lean{KummerCriterion.bernoulli_div_sModEq_of_modEq_full}
  \leanok
Let \(p\) be an odd prime.  Let \(m,n>0\) be even natural numbers such that
\[
  m \equiv n \pmod {p-1}
  \qquad\text{and}\qquad
  (p-1)\nmid n.
\]
Then there exists \(z\in \Z_p\) with
\[
  \frac{B_m}{m} - \frac{B_n}{n}
  =
  pz
  \qquad\text{in }\Q_p.
\]
\end{theorem}

\begin{proof}\leanok
The proof first reduces to \(p\ge 5\).  If \(p=3\), then \(p-1=2\) divides
every even exponent \(n\), contradicting the non-boundary hypothesis
\((p-1)\nmid n\).

For \(p\ge 5\), the proof is a strengthened elementary Voronoi argument.  The
first input is the strong Faulhaber congruence
\[
  \sum_{x=0}^{p-1} x^h - pB_h
  \in h p^2\Z_p
\]
for positive even \(h\) with \((p-1)\nmid h\).  In Lean this is
the strong Faulhaber lemma%
\lean{KummerCriterion.sum_range_pow_sub_p_mul_bernoulli_strong}.
It expands the power sum by mathlib's Faulhaber formula and proves that every
non-leading summand contains the required \(hp^2\)-factor.  The delicate
terms are handled with von Staudt-Clausen denominator control and binomial
divisibility.

The second input is the strong Voronoi congruence
\[
  (a^h-1)\frac{B_h}{h}
  -
  a^{h-1}\sum_{x=0}^{p-1} x^{h-1}\left\lfloor\frac{xa}{p}\right\rfloor
  \in p\Z_p
\]
for every \(a\) prime to \(p\), formalized as
the strong Voronoi theorem%
\lean{KummerCriterion.voronoi_congruence_mod_p_strong}.
This substitutes the strong Faulhaber congruence into a higher-modulus
Voronoi sum identity and cancels the nonzero rational factor \(hp\) in
\(\Q_p\).

Next choose a generator \(g\) of \((\Z/p\Z)^\times\), represented by a natural
number \(a\).  Since \(g\) has order \(p-1\), the condition
\((p-1)\nmid h\) implies \(a^h-1\) is a \(p\)-adic unit.  Dividing the strong
Voronoi congruence by this unit proves
\[
  \frac{B_h}{h}\in \Z_p,
\]
recorded as
the Voronoi integrality lemma%
\lean{KummerCriterion.bernoulli_div_self_mem_padicInt_of_not_sub_one_dvd_voronoi}.

Finally apply the strong Voronoi congruence to \(m\) and \(n\).  The factors
\(a^m\) and \(a^n\) agree because \(m\equiv n\pmod{p-1}\), and the floor sums
are congruent modulo \(p\) because their predecessor exponents are congruent;
this is the floor-sum comparison%
\lean{KummerCriterion.voronoi_floor_sum_sModEq_of_pred_modEq}.
The already-proved \(p\)-integrality of \(B_m/m\) and \(B_n/n\) lets the
comparison be carried out inside \(\Z_p\), and the difference is a multiple
of \(p\) in \(\Q_p\).
\end{proof}

\section{From divided Bernoulli numerators to irregularity}

The finite-set construction produces a prime divisor of the numerator of
\(B_M/M\), not directly a Kummer-range numerator \(B_{2k}\).  The next result
is the bridge from that unbounded witness to the usual definition of an
irregular prime.

\begin{theorem}[Carlitz divisor criterion]
  \label{thm:carlitz-divisor-criterion}
  \lean{KummerCriterion.not_isRegularPrime_iff_exists_dvd_bernoulli_div_self_num}
  \leanok
  \uses{thm:full-divided-kummer, thm:kummer-final}
Let \(q\) be an odd prime.  Then \(q\) is not regular if and only if there is
a positive even integer \(m\) such that
\[
  q \mid \operatorname{num}\left(\frac{B_m}{m}\right).
\]
\end{theorem}

\begin{proof}\leanok
If \(q\) is not regular, Kummer's criterion gives a Kummer-range index \(k\)
with
\[
  1\le k,\qquad 2k\le q-3,\qquad q\mid (B_{2k})_{\mathrm{num}}.
\]
Set \(m=2k\).  The range bound implies \(q\nmid m\), so \(m\) is a unit in
\(\Z_q\).  Dividing by this unit preserves numerator divisibility, giving
\[
  q\mid \operatorname{num}\left(\frac{B_m}{m}\right).
\]
The Lean helper is
the divisibility-through-division lemma%
\lean{KummerCriterion.dvd_bernoulli_div_self_num_of_dvd_bernoulli_num}.

Conversely assume \(q\mid\operatorname{num}(B_m/m)\) for some positive even
\(m\).  Von Staudt-Clausen excludes the boundary:
\[
  (q-1)\nmid m,
\]
because if \(q-1\mid m\), then \(q\) cannot divide the reduced numerator of
\(B_m/m\).  This is
the boundary-exclusion lemma%
\lean{KummerCriterion.sub_one_not_dvd_of_dvd_num_bernoulli_div_self}.
Let \(m'\) be the least nonnegative residue of \(m\) modulo \(q-1\).  The
boundary exclusion makes \(m'>0\), and parity gives that \(m'\) is even.  It
also satisfies \(m'<q-1\), hence
\[
  1\le m'/2,\qquad 2(m'/2)\le q-3.
\]

Apply \cref{thm:full-divided-kummer} to \(m\) and \(m'\).  Since
\(B_m/m\equiv B_{m'}/m'\pmod q\), the numerator divisibility transfers from
\(B_m/m\) to \(B_{m'}/m'\), and then to \(B_{m'}\).  The transfer is formalized
by
the \(p\)-adic numerator-transfer lemma%
\lean{KummerCriterion.dvd_bernoulli_num_of_padic_congruent_residue}.
Now \(m'=2(m'/2)\) lies in Kummer's range, so Kummer's criterion gives that
\(q\) is not regular.
\end{proof}

\section{The finite-set construction}

The escape argument starts from a finite set \(S\) and builds an even integer
which is divisible by \(q-1\) for every prime \(q\in S\).

\begin{definition}[Divisor-closed base]
  \label{def:irregular-base}
  \lean{KummerCriterion.irregularBase}
  \leanok
For a finite set \(S\subset\N\), set
\[
  C(S)=2\cdot \bigl(\max(3,\sup S)\bigr)!.
\]
\end{definition}

\begin{proof}\leanok
The factor \(2\) makes \(C(S)\) even, and the factorial makes
\[
  q-1\mid C(S)
\]
for every prime \(q\in S\).  The stronger closure property used internally is
that if a prime \(q\) divides \(C(S)\), then \(q-1\mid C(S)\) as well.
\end{proof}

\begin{theorem}[A new numerator prime outside a finite set]
  \label{thm:numerator-prime-for-carlitz-base}
  \lean{KummerCriterion.exists_numerator_prime_for_carlitz_base}
  \leanok
  \uses{def:irregular-base}
For every finite set \(S\subset\N\), there are natural numbers \(M\) and \(p\)
such that
\[
\begin{gathered}
  p\text{ is prime},\qquad p\ne 2,\qquad M>0,\qquad M\text{ is even},\\
  \forall q\in S,\ q\text{ prime}\Longrightarrow q-1\mid M,\\
  p\mid \operatorname{num}\left(\frac{B_M}{M}\right),
  \qquad p\notin S.
\end{gathered}
\]
\end{theorem}

\begin{proof}\leanok
Start with \(C=C(S)\).  The growth theorem
for divided Bernoulli numbers%
\lean{KummerCriterion.exists_large_even_multiple_abs_bernoulli_div_self_gt_one}
chooses \(t\) such that
\[
  \left|\frac{B_{C2^t}}{C2^t}\right|>1.
\]
Set \(M=C2^t\).  This \(M\) is positive and even, and it is still divisible by
\(q-1\) for every prime \(q\in S\).

Since the rational number \(B_M/M\) has absolute value greater than \(1\), its
reduced numerator has absolute value greater than \(1\), hence has a prime
divisor \(p\), by
the numerator-prime lemma%
\lean{KummerCriterion.exists_prime_dvd_num_of_one_lt_abs}.

It remains to prove that this prime was not already in \(S\).  Suppose
\(p\in S\).  Then \(p-1\mid M\) by construction.  But von Staudt-Clausen gives
the exclusion
\[
  p-1\mid M
  \quad\Longrightarrow\quad
  p\nmid \operatorname{num}\left(\frac{B_M}{M}\right),
\]
namely
the von Staudt numerator-exclusion lemma%
\lean{KummerCriterion.not_dvd_num_bernoulli_div_self_of_sub_one_dvd}.
This contradicts the choice of \(p\).  The same exclusion also proves
\(p\ne 2\), since \(p=2\) would make \(p-1=1\) divide \(M\).
\end{proof}

\section{Escaping every finite set}

\begin{theorem}[Finite-set escape]
  \label{thm:irregular-prime-outside-finite-set}
  \lean{KummerCriterion.exists_not_isRegularPrime_not_mem_carlitz}
  \leanok
  \uses{thm:numerator-prime-for-carlitz-base, thm:carlitz-divisor-criterion}
For every finite set \(S\subset\N\), there is a prime \(p\notin S\) such that
\(p\) is not regular.
\end{theorem}

\begin{proof}\leanok
Apply \cref{thm:numerator-prime-for-carlitz-base} to \(S\), obtaining a prime
\(p\notin S\) and a positive even \(M\) with
\[
  p\mid\operatorname{num}\left(\frac{B_M}{M}\right).
\]
The same theorem gives \(p\ne 2\).  Therefore the Carlitz divisor criterion
\cref{thm:carlitz-divisor-criterion} applies and proves that \(p\) is not
regular.
\end{proof}

\section{Infinitude}

\begin{theorem}[Infinitely many irregular primes]
  \label{thm:infinitely-many-irregular-primes}
  \lean{KummerCriterion.infinite_not_isRegularPrime}
  \leanok
  \uses{thm:irregular-prime-outside-finite-set}
The set
\[
  \{p:\N \mid p\text{ is prime and }p\text{ is not regular}\}
\]
is infinite.
\end{theorem}

\begin{proof}\leanok
Assume the set is finite, and let \(S\) be a finite set containing all of its
elements.  By \cref{thm:irregular-prime-outside-finite-set}, there is a prime
\(p\notin S\) which is not regular.  This contradicts the defining property of
\(S\).  Hence no finite set covers all irregular primes, so the set is
infinite.  The finite-set-to-infinite conversion is the elementary lemma
that no finite cover implies infinitude%
\lean{KummerCriterion.infinite_of_forall_finite_set_not_cover}.
\end{proof}

The checked proof contains no bundled Kummer provider or side-condition
package.  The only source bridge into irregularity is Kummer's criterion from
\cref{ch:kummer-final}, and the unbounded congruence used by Carlitz is the
public theorem \cref{thm:full-divided-kummer}.


\begin{thebibliography}{9}

\bibitem{carlitz1954}
L.~Carlitz,
\emph{Note on irregular primes},
Proc. Amer. Math. Soc. \textbf{5} (1954), 329--331.

\bibitem{ireland-rosen}
K.~Ireland and M.~Rosen,
\emph{A Classical Introduction to Modern Number Theory},
2nd ed., Graduate Texts in Mathematics 84, Springer-Verlag, New York, 1990.

\bibitem{jensen1915}
K.~L.~Jensen,
\emph{Om talteoretiske Egenskaber ved de Bernoulliske Tal},
Nyt Tidsskrift for Matematik \textbf{26} (1915), 73--83.

\bibitem{kummer1850}
E.~E.~Kummer,
\emph{Allgemeiner Beweis des Fermatschen Satzes, da{\ss} die Gleichung
$x^\lambda + y^\lambda = z^\lambda$ durch ganze Zahlen unl\"osbar ist,
f\"ur alle diejenigen Potenz-Exponenten $\lambda$, welche ungerade
Primzahlen sind und in den Z\"ahlern der ersten $(\lambda-3)/2$
Bernoullischen Zahlen als Factoren nicht vorkommen},
J. Reine Angew. Math. \textbf{40} (1850), 130--138.

\bibitem{sinnott1978}
W.~Sinnott,
\emph{On the Stickelberger ideal and the circular units of a cyclotomic
field},
Ann. of Math. (2) \textbf{108} (1978), 107--134.

\bibitem{washington}
L.~C.~Washington,
\emph{Introduction to Cyclotomic Fields},
2nd ed., Graduate Texts in Mathematics 83, Springer-Verlag, New York, 1997.

\end{thebibliography}
